% =============================================================================
%  Chapter 1 -- Introduction
% =============================================================================
%
%  LABELS. Every label carries a type prefix, so a reference reads unambiguously
%  and you can tell from \ref what will come back:
%      \label{sec:...}   chapters, sections, subsections
%      \label{fig:...}   figures
%      \label{tab:...}   tables
%      \label{eq:...}    equations
%      \label{alg:...}   algorithms (pseudo-code floats)
%      \label{lst:...}   listings (verbatim source)
%  Reference them with \Cref{sec:foo} ("Section 2") or \cref{} mid-sentence.
%  Never write "see section \ref{...}" by hand: cleveref (loaded by the class)
%  supplies the word, and keeps it correct if the float moves.
% =============================================================================

% The body restarts page numbering at 1 in arabic. This block belongs to the
% FIRST body chapter only -- do not copy it into chapters 2-5.
\blankpage
\setcounter{page}{1}
\pagenumbering{arabic}

\chapter{Introduction}
\label{sec:introduction}

% A template note. Delete it along with this comment once you start writing.
\emph{The number of chapters in this template, and the names they carry, are a
suggestion rather than a requirement. Rename them, merge them, split them or
drop them as the work demands. What follows is a common arrangement, not a prescribed one.}

% A second template note. Delete it too.
\emph{The same goes for spelling. British, Oxford and American English are all
acceptable; what is not acceptable is using two of them in one thesis, which is
the usual outcome of never having decided. Pick one before you start writing,
and set your spell checker to match. The template is set to British
via babel; if you choose American, change that option, since babel also decides
hyphenation and the words it generates for you.}

\section{Motivation}
\label{sec:motivation}

% STYLE: the introduction answers four questions, in this order, and each
% deserves roughly a paragraph or two. Resist the urge to start with a history
% of the field.
%
%   1. What is the problem, and why should anyone care?
%   2. What has been tried, and where specifically does it fall short?
%      (One paragraph. The detailed treatment goes in Chapter 2.)
%   3. What do you do about it? State the contribution plainly.
%   4. What did you find? Yes, give away the result here. A thesis is not a
%      detective novel.
%
% Write in the past tense for what you did ("we trained"), present tense for
% what is true ("the model generalises"). Pick "we" or the passive voice and
% hold it for the whole thesis; do not alternate.

Placeholder prose. Replace this section following the four-question structure
in the comments above.

\section{Research question}
\label{sec:research_question}

% STYLE: state the research question explicitly and set it off so a reader can
% find it without hunting. A short paragraph at most -- ideally the question
% itself in one sentence, with a sentence or two around it to fix scope and
% terms. If it runs longer than that, it is doing work that belongs in the
% motivation above. This is the spine of the thesis:
% every chapter that follows should be traceable back to it, and
% \Cref{sec:conclusion} has to answer it in as many words.
%
% A usable research question is:
%   * answerable with the evidence you can actually gather -- not "is X good?"
%     but "does X outperform Y on Z, and under which conditions?";
%   * narrow enough that a wrong answer is possible. If no result could have
%     contradicted it, it is a topic, not a question;
%   * free of the answer. "Why does X improve Y?" presumes it does -- establish
%     that first, or ask whether it does.
%
% Sub-questions are fine, and often clearer than one overloaded sentence: state
% the main question, then two or three that decompose it. Do not exceed that --
% a thesis that asks five questions usually answers none of them.

This thesis addresses the following research question:
\begin{quote}
  \emph{One sentence, ending in a question mark, that the rest of the thesis
  exists to answer. Where the question needs it, one or two further sentences
  fixing its scope: the setting it applies to, or what a term in it is taken to
  mean. No more than a short paragraph in total.}
\end{quote}
% Optional decomposition -- delete if the question above stands on its own.
It is approached through two sub-questions:
\begin{enumerate}[itemsep=2pt,topsep=4pt]
  \item A narrower question whose answer the main question depends on.
  \item A second one.
\end{enumerate}

\section{Contributions}
\label{sec:contributions}

% STYLE: an explicit, numbered list of contributions is the single highest-value
% paragraph in the thesis -- it is what an examiner looks for first. Each item
% should be something you could defend as novel or as real work performed.

This thesis makes the following contributions:
\begin{enumerate}[itemsep=2pt,topsep=4pt]
  \item A concrete, verifiable claim about something you built or showed.
  \item A second one. Three to five items is the usual range.
  \item A third one.
\end{enumerate}

\section{Outline}
\label{sec:outline}

% STYLE: one sentence per chapter, naming what the reader will get from it.
% Keep it short -- this section is a signpost, not a summary.

The front matter carries four reference lists, and they divide into two kinds.
The \emph{List of Figures} and \emph{List of Tables} are indexes: they let a
reader who remembers a diagram find it again, and they are generated
automatically from the short form of each caption, so they need no maintenance
beyond writing good captions. The \emph{List of Notations} and \emph{List of
Abbreviations} are glossaries: they let a reader who has forgotten what
$\mathcal{L}$ or \textsc{ner} means look it up without hunting back through
the text for its first use. Those two are written by hand in
\texttt{text/notations.tex} and \texttt{text/abbreviations.tex}, and nothing
checks them, so a symbol renamed in a later chapter goes stale in the list
silently. Keep them in step as you write rather than at the end, and delete
either list if the thesis has no notations or abbreviations. Neither
replaces introducing a symbol or expanding an abbreviation at first use in the
prose.

\Cref{sec:background} reviews prior work on \dots{}
\Cref{sec:methodology} introduces \dots{}
\Cref{sec:experiments} evaluates \dots{}
\Cref{sec:conclusion} concludes and discusses limitations.
